\newpage
\phantomsection
\label{prf:topological-basis-pointwise-characterization}
\LRAProofFor{prop:topological-basis-pointwise-characterization}

\begin{remark*}[Return]
\hyperref[prop:topological-basis-pointwise-characterization]{Return to Proposition}
\end{remark*}

\begin{proposition*}[Pointwise Characterization of a Basis]
Let \((X,\tau)\) be a topological space, and let
\(\mathcal{B}\subseteq\tau\).  Then \(\mathcal{B}\) is a basis for \(\tau\)
if and only if for every open set \(O\in\tau\) and every \(x\in O\), there is
\(B\in\mathcal{B}\) such that \(x\in B\subseteq O\).  Consequently, a subset
\(A\subseteq X\) is open if and only if for each \(x\in A\), there is
\(B\in\mathcal{B}\) such that \(x\in B\subseteq A\).
\end{proposition*}

\begin{proof}[Professional Standard Proof]
\LRAProofBodyStart
Suppose first that \(\mathcal{B}\) is a basis for \(\tau\).  Let
\(O\in\tau\) and let \(x\in O\).  Since \(\mathcal{B}\) is a basis, there is
a subcollection \(\mathcal{A}\subseteq\mathcal{B}\) such that
\[
  O=\bigcup_{B\in\mathcal{A}}B.
\]
Because \(x\in O\), there is some \(B\in\mathcal{A}\) with \(x\in B\).  Since
\(\mathcal{A}\subseteq\mathcal{B}\) and its union is \(O\), this set satisfies
\(B\in\mathcal{B}\) and \(x\in B\subseteq O\).

Conversely, assume the pointwise condition.  Let \(O\in\tau\).  For each
\(x\in O\), choose \(B_x\in\mathcal{B}\) such that \(x\in B_x\subseteq O\).
Then
\[
  O=\bigcup_{x\in O}B_x,
\]
with the empty union understood when \(O=\emptyset\).  Hence every open set is
a union of elements of \(\mathcal{B}\), so \(\mathcal{B}\) is a basis for
\(\tau\).

For the consequence, if \(A\in\tau\), apply the pointwise condition with
\(O=A\).  Conversely, if each \(x\in A\) has a basis element
\(B_x\in\mathcal{B}\) with \(x\in B_x\subseteq A\), then
\[
  A=\bigcup_{x\in A}B_x.
\]
Each \(B_x\) is open because \(\mathcal{B}\subseteq\tau\), and arbitrary
unions of open sets are open.  Therefore \(A\in\tau\).
\end{proof}

\begin{proof}[Detailed Learning Proof]
\LRAProofBodyStart
A basis builds open sets by unions.  If an open set \(O\) is already written
as a union of basis elements, then any point \(x\in O\) must lie in one of
those basis elements.  That basis element is automatically contained in
\(O\), because it is one of the pieces used in the union.

For the reverse direction, assume that every point of every open set has a
basic open set around it that stays inside the open set.  Collect one such
basis element \(B_x\) for each point \(x\in O\).  Their union contains every
point of \(O\), and none of them leaves \(O\).  Therefore their union is
exactly \(O\).  This proves that every open set is a union of basis elements.

The same argument characterizes openness of an arbitrary subset \(A\).  If
\(A\) is open, the pointwise basis condition applies to \(A\).  If every point
of \(A\) lies in a basis element contained in \(A\), then \(A\) is a union of
open basis elements and is therefore open.
\end{proof}

\begin{remark*}[Proof structure]
Use the union definition of a basis in one direction and construct the union
of all local basis elements in the other direction.
\end{remark*}

\begin{dependencies}
\begin{itemize}
  \item \hyperref[def:topological-space]{Definition: Topological Space}
  \item \hyperref[def:topological-basis]{Definition: Basis for a Topology}
\end{itemize}
\end{dependencies}

\clearpage
